\documentclass[11pt]{article}

\usepackage{amsmath,amssymb,amsthm,mathtools}
\usepackage[T1]{fontenc}
\usepackage{lmodern}
\usepackage{geometry}
\usepackage{enumitem}
\usepackage{hyperref}
\usepackage{microtype}
\geometry{margin=1in}
\hypersetup{colorlinks=true, linkcolor=blue, citecolor=blue, urlcolor=blue}

\newtheorem{theorem}{Theorem}
\newtheorem{proposition}[theorem]{Proposition}
\newtheorem{corollary}[theorem]{Corollary}
\newtheorem{remark}[theorem]{Remark}

\newcommand{\Zm}{\mathbb Z_m}

\title{A cleanup note for the odd-$m$ D5 globalization theorem}
\author{}
\date{March 2026}

\begin{document}
\maketitle

\begin{abstract}
We record the current accepted theorem-level form of the odd-$m$ $d=5$
package. The final claim is that raw global $(\beta,\delta)$ is exact on the
true accessible boundary union, equivalently that every actual lift of the
exceptional cutoff row $\delta=3m-3$ continues through
$3m-2\to 3m-1$ and lands in the regular continuing endpoint class there.
The note is written as a cleanup target: it keeps the structural input,
componentwise bridge, globalization reductions, and final gluing theorem
separate, and it preserves the scope distinction that the bridge is
componentwise before the globalization argument closes.
\end{abstract}

\section{Statement and scope}

Throughout, fix an odd modulus $m$ in the accepted D5 regime.
We write the abstract right-congruence coordinate as $\rho$ and the concrete
boundary coordinate as $\delta$.
The purpose of this note is only to record the final theorem chain cleanly.
It does not attempt to retell the full historical development of the D5 branch.

The final theorem-level claim is:
\begin{quote}
raw global $(\beta,\delta)$ is exact on the true accessible boundary union.
\end{quote}
Equivalently:
\begin{enumerate}[label=\textup{(\roman*)},leftmargin=2.5em]
\item the abstract bridge $(\beta,\rho)$ is exact on the true accessible
boundary union;
\item the concrete bridge $(\beta,\delta)$ is exact globally, not merely
componentwise;
\item $\rho=\rho(\delta)$ globally;
\item there is no hidden second endpoint sheet over $\delta=3m-1$.
\end{enumerate}

The sharp row formulation is:
\begin{quote}
Every actual lift of the exceptional cutoff row $\delta=3m-3$ continues through
$3m-2\to 3m-1$ and lies in the regular continuing endpoint class there.
\end{quote}

\section{Imported structural inputs}

The final theorem uses the following structural inputs.

\begin{proposition}[explicit trigger family]
On the unresolved best-seed channel $R1\to H_{L1}$, the target hole family is
\[
H_{L1}=\{(q,w,u,\lambda)=(m-1,m-1,u,2):u\neq 2\}.
\]
\end{proposition}

\begin{proposition}[universal first-exit targets]
The active best-seed branch first exits to $H_{L1}$ at exactly one of the two
states
\[
T_{\mathrm{reg}}=(m-1,m-2,1),
\qquad
T_{\mathrm{exc}}=(m-2,m-1,1).
\]
Regular source families exit at $T_{\mathrm{reg}}$ by direction $2$, while the
exceptional source family exits at $T_{\mathrm{exc}}$ by direction $1$.
\end{proposition}

The role of these propositions in the final proof is narrow but decisive:
first-exit data are used only to certify the exceptional continuation through
the interface row.

\section{The concrete bridge package}

On the $\Theta=2$ boundary section define
\[
\sigma \equiv w+u-q-1 \pmod m,
\qquad
\delta := q+m\sigma.
\]

\begin{proposition}[componentwise concrete bridge]
On each splice-connected accessible component:
\begin{enumerate}[label=\textup{(\alph*)},leftmargin=2.5em]
\item the odometer splice law holds,
\[
\delta'=\delta+1 \pmod{m^2};
\]
\item the current event class $\epsilon_4$ is a function of $(\beta,\delta)$.
\end{enumerate}
In particular, $(\beta,\delta)$ is an exact concrete bridge coordinate
componentwise.
\end{proposition}

This bridge proposition is intentionally stated only componentwise. The global
upgrade comes later.

\section{Globalization reductions}

\begin{proposition}[tail-length reduction]
Once the concrete bridge package is accepted, two realizations of the same
boundary label $\delta$ can differ only by remaining full-chain tail length.
Equivalently, there is no further local future-law ambiguity at fixed realized
$\delta$.
\end{proposition}

\begin{proposition}[regular continuation]
Every regular realization of every realized $\delta$ continues on the true
larger regular accessible union.
\end{proposition}

\begin{proposition}[exceptional interface theorem]
At chart/interface level, the exceptional continuation is pinned to
\[
3m-3 \to 3m-2 \to 3m-1.
\]
The labels $3m-2$ and $3m-1$ are the distinguished regular interface labels:
$3m-2$ is the terminal regular source-$4$ occurrence and $3m-1$ is the regular
source-$1$ start.
\end{proposition}

\begin{proposition}[single-row reduction]
The full raw globalization theorem is equivalent to one exceptional-row
statement: every actual lift of the exceptional cutoff row $\delta=3m-3$ must
glue through $3m-2\to 3m-1$ into the regular continuing class.
\end{proposition}

These reductions isolate the only remaining obstruction after the regular
sector is closed: the exceptional cutoff row cannot be allowed to launch a new
endpoint sheet.

\section{Final theorem}

\begin{theorem}[odd-$m$ D5 final globalization theorem]
For odd $m$ in the accepted D5 regime, every actual lift of the exceptional
cutoff row $\delta=3m-3$ continues through
\[
3m-2 \to 3m-1
\]
and lies in the regular continuing endpoint class there.
Consequently:
\begin{enumerate}[label=\textup{(\arabic*)},leftmargin=2.5em]
\item there is no mixed-status realized $\delta$;
\item fixed realized $\delta$ determines tail length;
\item fixed realized $\delta$ determines endpoint class and future word;
\item $\rho=\rho(\delta)$ globally;
\item raw global $(\beta,\delta)$ is exact.
\end{enumerate}
\end{theorem}

\begin{proof}
By regular continuation, no regular full chain is terminal.
By the componentwise splice law, every exceptional full chain before the cutoff
continues internally.
So only the exceptional cutoff chain at $\delta=3m-3$ needs work.

The universal first-exit theorem gives a unique actual exceptional exit target
and branch direction on the active branch. The interface theorem then pins the
corresponding continuation in chain labels to
\[
3m-3 \to 3m-2 \to 3m-1.
\]
Hence the exceptional cutoff chain also has a forward successor.
Therefore every accessible full chain has a forward successor.

If some splice-connected accessible component were finite, its realized
boundary image in $\mathbb Z/m^2\mathbb Z$ would have a right-end chain with no
successor, contradicting the previous paragraph. Hence every splice-connected
accessible component is total, and every realization has infinite remaining
tail length.

Now apply the tail-length reduction: at fixed realized $\delta$, the only
possible ambiguity was tail length. Since tail length is now always infinite,
fixed realized $\delta$ determines the endpoint class and future word.
Therefore $\rho=\rho(\delta)$ globally, which is exactly the claim that raw
global $(\beta,\delta)$ is exact.
\end{proof}

\section{Concrete $m=5$ specialization}

The theorem has a directly checkable row form at $m=5$.
Since
\[
3m-3=12,
\qquad
3m-2=13,
\qquad
3m-1=14,
\]
the exceptional interface statement becomes:
\begin{quote}
Every actual lift of $\delta=12$ continues through $13\to 14$ and glues into
the regular continuing endpoint class at $14$.
\end{quote}

Using
\[
\delta=q+5\sigma,
\]
this is the row
\[
12=2+5\cdot 2,
\qquad
13=3+5\cdot 2,
\qquad
14=4+5\cdot 2,
\]
so the interface lies on the slice $\sigma=2$ with $q=2,3,4$.
The corresponding trigger family and first-exit targets are
\[
H_{L1}=\{(4,4,u,2):u\neq 2\},
\qquad
T_{\mathrm{reg}}=(4,3,1),
\qquad
T_{\mathrm{exc}}=(3,4,1).
\]

\begin{remark}[separate explicit witness at $m=5$]
The bundle also contains a separate concrete file
\texttt{d5\_m5\_kempe\_witness\_26.json} recording a 26-move Kempe-valid witness
with cycle counts $(1,1,1,1,1)$ on $m=5$.
That object is a direct small-modulus existence witness.
It is useful supporting evidence, but it is logically distinct from the final
raw-state globalization theorem proved above.
\end{remark}

\section{Cleanup warning}

In any later manuscript integration, the following scope distinctions should
remain explicit.

\begin{itemize}[leftmargin=2em]
\item The concrete bridge package is componentwise before globalization.
\item The exceptional interface theorem is chart/interface level before it is
combined with the actual first-exit theorem.
\item The global theorem is obtained only after combining the structural
first-exit input, the interface theorem, the regular continuation theorem, and
the tail-length reduction.
\end{itemize}

\end{document}
